\section{Клиенты}
\label{sec:clients}

\textbf{Весовые категории.} Гибрид имеет смысл для моделей, которые вместе с кэшем нескольких сессий помещаются в
память персональной машины. Модели в триллион параметров и больше на такие машины не помещаются и остаются
облачными, а маленькие модели клиент читает достаточно быстро и сам. Между ними три весовые категории, и каждая
совпадает с распространённым объёмом быстрой унифицированной памяти (таблица~\ref{tab:tiers}): плотные модели около
30 миллиардов параметров для 64~ГБ, модели MoE около 180 миллиардов для 128~ГБ и около 330 миллиардов для 256~ГБ. Это
массовый класс гибридного инференса: машину такого класса может позволить себе отдельный специалист или небольшая
команда.

\begin{table}[t]
\centering
\caption{Весовые категории гибрида: веса в 4 битах и число сессий, которые помещаются рядом с ними, если 12~ГБ
оставить системе. Для Qwen3.8-27B миллион токенов лежит за пределами штатного контекста, это оценка по памяти.}
\label{tab:tiers}
\footnotesize
\begin{tabular}{llrrrr}
\toprule
Память & Модель & Веса, ГБ & Сессий по 256 тыс. & По 500 тыс. & По 1 млн \\
\midrule
64 ГБ & Qwen3.8-27B, плотная & 19 & 3 & 1 & --- \\
128 ГБ & Qwen3.8-27B, плотная & 19 & 10 & 5 & 2 \\
128 ГБ & Qwen3.8-Flash-Next, MoE & 70 & 13 & 7 & 3 \\
256 ГБ & GLM-5.3-Flash, MoE & 181 & 37 & 19 & 10 \\
\bottomrule
\end{tabular}
\end{table}

Больше всего схема даёт машине, у которой высокая полоса памяти и слабые вычислители. Так работает Qwen3.8-27B на M3
и M4 Max: модель помещается, генерация идёт с приемлемой скоростью, а чтение глубокого контекста остаётся узким
местом, которое схема переносит на сервер.

\textbf{Память под контекст.} Веса Qwen3.8-27B занимают около 19~ГБ в 4-битной квантизации и около 15~ГБ в 3-битной.
Каждая тысяча токенов контекста добавляет 35~МБ кэша внимания, и к каждой сессии добавляется 157~МБ рекуррентного
состояния. Таблица~\ref{tab:fit} показывает, сколько сессий помещается в память машины, если 12~ГБ оставить
системе и приложениям. На 36~ГБ помещается одна сессия на 128 тысяч токенов, на 64~ГБ уже семь таких сессий или одна на
полмиллиона, на 128~ГБ по объёму памяти две сессии на миллион. Для Qwen3.8-27B столбцы 500 тысяч и миллион
остаются оценкой по памяти: штатный контекст модели 262 тысячи, а более длинный требует расширения контекста и
проверки качества (таких проверок мы не делали).

\begin{table}[t]
\centering
\caption{Сессий Qwen3.8-27B в памяти клиента: веса в 4 битах (около 19~ГБ), кэш внимания в 8 битах (34\,844 байта на
токен), рекуррентное состояние 157~МБ на сессию, 12~ГБ оставлено системе и приложениям.}
\label{tab:fit}
\footnotesize
\begin{tabular}{lrrrrrr}
\toprule
Память & 32 тыс. & 64 тыс. & 128 тыс. & 256 тыс. & 500 тыс. & 1 млн \\
\midrule
36 ГБ & 3 & 2 & 1 & --- & --- & --- \\
64 ГБ & 25 & 13 & 7 & 3 & 1 & --- \\
96 ГБ & 51 & 27 & 14 & 7 & 3 & 1 \\
128 ГБ & 76 & 40 & 21 & 10 & 5 & 2 \\
\bottomrule
\end{tabular}
\end{table}

\textbf{Скорость генерации и её запас.} На каждом шаге генерация читает веса, которые участвуют в шаге, и весь кэш
контекста, поэтому её предел задаёт полоса памяти. У Qwen3.8-27B это все веса, около 19~ГБ. На M3 Max (полоса
400~ГБ/с) наш замер дал 7.5--8.8 токена в секунду на глубине 16 тысяч, около 40\% от предела, который получается
делением полосы на объём читаемых за шаг данных. Несколько потоков читают веса один раз на всех, поэтому суммарный
предел растёт с числом потоков: для четырёх потоков на глубине 32 тысячи он около 64 токенов в секунду. Замер дал 13.2 токена в секунду по времени запросов (14.05 по интервалам генерации) против 8.44 у одного
потока, около 20\% этого предела. Насколько движки сократят разрыв между
замером и пределом, покажут следующие замеры; расчёты в этой работе используют замеренную скорость. Когда чтение
порций уходит на сервер, работа клиента сводится к генерации, и каждый выигрыш движка в генерации становится
выигрышем всей схемы.

\textbf{Сколько успевает клиент за день.} Таблица~\ref{tab:clients} показывает, сколько сессий нашего профиля клиент
проводит за восьмичасовой день с четырьмя параллельными агентами, сам и в гибриде. При нескольких агентах большую
часть времени клиента занимает генерация: ответ в 512 токенов при четырёх потоках на M3 Max генерируется около двух с половиной минут даже на глубине 32 тысячи, а чтение порции в среднем занимает 28 секунд. Поэтому за день гибрид ускоряет работу M3 Max в 1.6
раза, а каждый отдельный агент ждёт prefill на 19--27 секунд меньше на вызов (таблица~\ref{tab:offload}). В GPU
поколения M5 есть блоки, которые ускоряют матричные вычисления; по нашей оценке такие клиенты читают контекст в
четыре раза быстрее M3 Max, и гибрид прибавляет им за день 11--17\%. Эту оценку нужно проверить замером. Mac Studio с M5
Ultra вдвое быстрее M5 Max и в чтении, и в генерации, поэтому соотношение между ними прежнее и выигрыш почти такой
же. NVIDIA DGX Spark на GB10 Grace Blackwell служит клиентом на CUDA со 128~ГБ унифицированной памяти LPDDR5X \cite{nvidiadgxspark}. По нашей оценке DGX Spark читает контекст быстро, а генерирует медленно из-за полосы 273~ГБ/с, на треть меньшей, чем у M3 Max, и выигрыш у него меньше всех; замеров на нём мы не делали. С 16 октября 2026 года продаются и ноутбуки RTX Spark на чипе N1X \cite{nvidiartxsparknews} с унифицированной памятью LPDDR5X на 256-битной шине, до 128~ГБ в старшей конфигурации \cite{nvidiartxspark}. Полосу памяти RTX Spark NVIDIA не публикует; RTX Spark мы не измеряли и в расчёты не включали. Модели
MoE с малым числом активных параметров, например Qwen3.8-Flash-Next с 6 миллиардами активных из 180, клиент и сам
читает достаточно быстро; к тому же им нужно 96--128~ГБ памяти.

\begin{table}[t]
\centering
\caption{Условных сессий по 116 вызовов (среднее число вызовов в сессии в журналах) за восьмичасовой день с четырьмя
параллельными агентами, Qwen3.8-27B, канал 1~Гбит/с.
Генерация одного потока на глубине 32 тысячи. Для M3 Max ёмкость рассчитана по замеренным кривым; для остальных клиентов prefill и генерация пересчитаны по
вычислительной мощности и полосе памяти. Все значения в таблице оценки.}
\label{tab:clients}
\footnotesize
\begin{tabular}{lrrrr}
\toprule
Клиент & Генерация, ток/с & Сессий сам & Сессий в гибриде & Выигрыш \\
\midrule
M3 Max 128 ГБ & 8.4 & 3.37 & 5.31 & 1.58 \\
M4 Max 128 ГБ & 11.5 & 5.24 & 7.18 & 1.37 \\
M5 Max 64 ГБ & 13.0 & 6.58 & 7.61 & 1.16 \\
M5 Max 128 ГБ & 13.0 & 6.86 & 8.06 & 1.17 \\
M5 Ultra 96 ГБ & 25.9 & 13.41 & 15.34 & 1.14 \\
M5 Ultra 256 ГБ & 25.9 & 14.13 & 15.64 & 1.11 \\
AMD Strix Halo 128 ГБ & 5.4 & 3.04 & 3.44 & 1.13 \\
NVIDIA DGX Spark & 5.8 & 3.52 & 3.66 & 1.04 \\
\bottomrule
\end{tabular}
\end{table}

Выходит, что на нашем профиле с несколькими агентами гибрид в первую очередь сокращает ожидание каждого
агента, а пропускную способность клиента за день поднимает умеренно. Главный выигрыш приходится на глубокий
контекст, где чтение на клиенте растягивается на десятки минут (раздел~\ref{sec:million}), и на машины, которые читают
контекст медленно. Машины класса Max и Ultra с плотной моделью работают в обоих режимах: они уже есть у многих
пользователей, без сервера работают автономно, а с сервером почти не ждут prefill.
